% Sections 4 to 6.

\section{The Life Cycle of a Transaction}\label{sec:tx}

This section describes what lives on the network and how payments work,
step by step, from the composition of a payment in the wallet to its
irreversibility. A private transaction worthy of the name must prove four
things without revealing anything: that the sender owned the funds, that
the amounts balance, that the coin being spent was never spent elsewhere,
and that the sender authorized the spending. Each proof rests on a
distinct cryptographic construction, inherited from the CryptoNote lineage
that Monero and NERVA have carried for a decade and validated under
continuous attack. The merit of that lineage is not the invention of
exotic objects: it is the composition of simple, individually proven
objects into a system whose privacy property is demonstrable rather than
declared.

\subsection{Step 1: the one-time address}

Every payment identity publishes two public keys, a viewing key $A$ and a
spending key $B$. When Alice pays Bob, she draws a random value $r$,
computes the point $R = rG$, then derives a one-time destination address
$P$ that nobody can link to Bob's published address. Bob scans every new
output of the DAG and recognizes his own by replaying the same computation
with his private viewing key $a$: one multiplication and one addition
suffice, and a wallet can scan thousands of outputs per second. Each
output also carries a one-byte view tag, the top byte of the shared
secret: for the outputs whose tag does not match, Bob's wallet skips the
curve derivation entirely, so scanning stays linear and cheap at a
two-second cadence. No address is ever reused; graph analysis, the first
weapon of transparent chains, has no node to hold on to.
\begin{equation}\label{eq:onetim}
P = H_s(rA)\,G + B
\qquad\text{and, on the recipient side,}\qquad
P = H_s(aR)\,G + B
\end{equation}
with $R = rG$ published in the transaction, $H_s$ a hash mapped onto the
curve, and $G$ the group generator.

\subsection{Step 2: amounts under commitment}

The amount $v$ of each output is never written in clear. It is sealed in
a Pedersen commitment: an envelope with a remarkable property, addition.
Two envelopes can be summed, and their sum commits to the sum of the
amounts, without anyone opening the envelopes. That homomorphy lets the
network verify conservation of amounts at every spend while remaining
unable to read a single amount. The blinding factor $b$, specific to each
output, forbids the brute-force guessing of small amounts, the historical
weakness of the first masked currencies; it also ties the output to its
recipient, who alone can reconstruct its value.
\begin{equation}\label{eq:pedersen}
C = vH + bG
\end{equation}

\subsection{Step 3: the ring of decoys}

A spend references prior outputs; referencing them directly would reveal
the graph. The ring signature therefore selects, around the real output
being spent, fifteen decoys drawn from the outputs already ordered by the
DAG, never from its unordered frontier (Section~\ref{sec:settlement}):
sixteen candidate keys, one of them real, and a signature proving
knowledge of one secret among the sixteen without saying which. The
probability of identifying the real spend is one in sixteen for an outside
observer, whatever the computing power; anonymity is probabilistic per
set, exactly as irreversibility is probabilistic per depth in classical
proof of work. Sixteen equals the strongest ring the lineage deploys in
production today, Monero's enforced size since 2022; it is a launch floor
rather than a ceiling, and the lineage's own research toward full-chain
membership proofs is the direction the network watches, ready to inherit
once it is battle-tested by others first. Decoys are sampled by age
brackets following the real distribution of spending, so that the
composition of the ring does not statistically betray the age of the real
output.

\subsection{Step 4: the key image, guard against double spending}

Without a guard, the same secret could sign two different rings and spend
the same output twice. The key image resolves that paradox: a
deterministic fingerprint of the spent output, published with the
transaction and stored in a nullifier registry. Two spends of the same
secret produce the same key image, the network rejects the second, and yet
the image allows neither tracing back to the real output of the ring nor
linking two payments of the same wallet: the construction is untraceable
and unlinkable. Double spending is thereby made structurally impossible
without sacrificing privacy, a property the CryptoNote lineage has
demonstrated in production since 2014.
\begin{equation}\label{eq:keyimage}
I = x\,H_p(xG)
\end{equation}
published and checked against the nullifier registry.

\subsection{Step 5: the range proof}

A commitment masks the amount, but nothing would forbid committing a
negative amount: spend ten units by fabricating an output of minus ten,
and the homomorphic arithmetic would still verify while creating money.
The range proof closes that door: each output proves, without revealing
$v$, that the amount belongs to the bounded positive interval of
authorized amounts. ANTUMBRA retains Bulletproofs+, the lineage standard
in production since 2022, whose size grows as the logarithm of the
precision and whose verification, aggregated per block, parallelizes on
commodity processors. A classical range proof would ruin the throughput;
the logarithmic form preserves it.
\begin{equation}\label{eq:range}
\text{proof: } v \in [0, 2^{64}) \text{ without revealing } v
\end{equation}

\subsection{Step 6: conservation, signature and fees}

The assembled transaction carries inputs (prior outputs spent in a ring),
outputs (fresh commitments), key images, range proofs, a bounded free-data
field, and fees. The central validation rule is one equality: the sum of
the input commitments, whose composition the sender can open, equals the
sum of the output commitments plus the fee commitment. Because commitments
add up, this equality proves conservation of the money mass without ever
unveiling an amount. The ring signature proves authorization, the key
images prove non-spending, the range proofs bound the amounts: the four
requirements of the section opening are each covered by a distinct
construction, and that separation is what makes the whole auditable
clause by clause.
\begin{equation}\label{eq:conservation}
\sum_{\mathrm{in}} C_j \;=\; \sum_{\mathrm{out}} C_i \;+\; C_{\mathrm{fees}}
\end{equation}

\subsection{Step 7: stem propagation}

Diffusing a transaction from the sender's IP address would hand the
network observer what the cryptography has just refused. Dandelion++
propagation routes the transaction through a stem phase, peer to peer on
a random path, before blooming into a fluff phase toward the whole
network: the apparent origin of the diffusion is no longer the sender but
an anonymous peer several hops away. Tor transport, enabled by default,
finishes the work by masking the network address of the nodes themselves.
An observer wanting to link a payment to a machine would have to defeat
simultaneously the ring, the commitment, the one-time address, the stem
and the anonymity network: five independent walls, each individually
proven for years.

\subsection{Step 8: inclusion, checkpoint and irreversibility}

The transaction enters a two-second block, produced by a CPU miner and
ordered by the DAG. The miner verifies the entirety of the proofs
described above before including the transaction, and every node replays
those verifications on reception: the network trusts nobody, it replays
everything. Inclusion is not yet irreversibility; that comes from the next
checkpoint, signed every four seconds by the Ring
(Section~\ref{sec:ring}). As soon as a checkpoint reaches its quorum,
everything it covers is finalized: a double spend would have to
simultaneously reorganize the accumulated work of the DAG and break the
majority of the committee's signatures. For the merchant at the counter,
the complete sequence, from sending to irreversibility, fits in under six
seconds.

\begin{table}[htbp]
\centering
\small
\caption{The transaction types of the network.}
\label{tab:txtypes}
\begin{tabularx}{\textwidth}{@{}L{3.6cm}YY@{}}
\toprule
Transaction type & Specific content & Additional validation \\
\midrule
Transfer & standard Veil outputs & none beyond the eight steps \\
Mandate & output with a release predicate (Section~\ref{sec:contracts}) & the predicate is executable and bounded in size \\
Identity registration & Ember or Cipher, sponsor, perimeter & valid sponsor, sponsorship quota, presence proof \\
Kleos attestation & testimony of $\pm 1$ with on-chain proof & established witness, Echo budget, anti-reciprocity \\
Governance vote & chamber, proposal, choice & voting right of the relevant chamber \\
Checkpoint & hash of the ordered DAG tip & quorum of 37 signatures out of 55 seats \\
\bottomrule
\end{tabularx}
\end{table}

\subsection{The fee model}

Fees purchase two real resources: room in blocks, which is bandwidth, and
verification work, which is processor time of the nodes. The retained
model is deliberately simple, because a simple model is predictable for
the merchant, auditable for the node, and hard to game: a fixed part, a
part proportional to size, and a surcharge per ringed output, the most
expensive construction to verify. Fees go entirely to miners; no protocol
levy is added, and the ordering of transactions in the DAG belongs to no
value extractor: the very structure of the DAG, where parallel blocks are
ordered by rule consensus, closes the front-running niche that funds the
mempool wars of builder chains.
\begin{equation}\label{eq:fees}
f = f_0 + f_{\mathrm{kB}} \cdot n_{\mathrm{kB}} + f_{\mathrm{ring}} \cdot m
\end{equation}
for $m$ ringed outputs and $n_{\mathrm{kB}}$ kibibytes; governance can
adjust $f_0$ within $[f_{\min}, f_{\max}]$ by the fast track, and outside
those bounds only by the constitutional track.

To close the section honestly: seven of the eight steps are inherited
directly from the CryptoNote lineage, proven in production on networks
processing hundreds of thousands of transactions per day. The
novelty of ANTUMBRA is the composition, the ring widened to sixteen, the
two-second cadence that the DAG sustains, and the coupling with the
finality of the Ring. The only genuinely new construction, checkpoints
anchored on reputation, is the subject of the next section and of a
dedicated risk treatment. That distribution of novelty, maximal on rules
and minimal on primitives, is a decision of method: new cryptographic
primitives are the graveyard of private currency projects, and the
roadmap introduces none before external audit.

\section{The Ring: Reputation-Anchored Finality}\label{sec:ring}

Being included in a two-second block is not enough: the payment must also
become irreversible fast enough for the customer standing at the counter.
Fault-tolerant-consensus chains solve that problem in exchange for a
capital lock: one must stake to sign, and whoever can stake buys
finality. The Ring takes the opposite path. Its fifty-five seats are not
for sale, they are drawn, at each era, among the identities whose Kleos
exceeds the candidacy threshold: a deterministic draw, weighted by score,
over a candidate window published in advance, auditable by everyone. The
draw formula is linear in score, deliberately: quadratic weighting would
concentrate seats at the top, and the concentration index of total Kleos,
monitored by governance, bounds any single actor's share of a cohort in
any case.
\begin{equation}\label{eq:draw}
P(i) = K_i \,\big/\, \sum_{j \in C} K_j
\end{equation}
the probability of drawing a seat, for candidates $C$ with Kleos of at
least 70.

Every four seconds, the Ring collectively signs a checkpoint, the hash of
the ordered DAG up to a given tip. The Byzantine fault tolerance of a
committee of fifty-five members admits eighteen faulters; the signature
quorum is therefore thirty-seven, the strict two thirds. As soon as a
checkpoint reaches quorum, everything it covers is finalized: a double
spend would have to simultaneously break the majority of the Ring and the
continuity of mining, two independent fault economies. A seat that signs
a competing fork is stripped: its Kleos reset to zero, and with it the
years that constituted it. The sanction does not fall on a recoverable
deposit; it falls on the only thing an ANTUMBRA validator owns that is
precious.
\begin{equation}\label{eq:bft}
n = 3f + 1 = 55,\quad f = 18,\quad \text{quorum} = 2f + 1 = 37 \text{ signatures}
\end{equation}

Failure is designed, not denied. If a third of the seats fall silent,
through outage, censorship or attack, the network does not stop: mining
continues, payments continue, and finality falls back on classical
proof-of-work depth, ten blocks, twenty seconds, until the next era
removes the silent seats. The Ring is not a production bottleneck, it
produces nothing: it is a notary of order, and a notary on strike slows
irreversibility, never payments. That decoupling is the reason an
egalitarian network can afford world-class finality without ever selling
its seats.

It must be said once and for all: this assembly is not proof of stake. No
capital is locked, no yield is paid to seats, no hierarchy of enrichment
reproduces itself. The power of finality derives from Kleos; Kleos
derives from behavior and time; and time is for sale nowhere. That is
the exact correction of the formula that inspired this project, where
reputation multiplied capital: here, capital multiplies nothing.

\subsection{The Genesis Ring: bootstrapping finality without privileges}

The strictness of R1 raises an honest question: if candidacy demands
fifteen incident-free eras, who signs the checkpoints during the seven
and a half years in which nobody can honestly qualify? The answer must
not be a founder key, an appointed council, or a lowered guard: each of
those would import, on day one, exactly the privilege the network
promises to abolish. The answer is a published ratchet, stricter than
the steady state, never weaker.

\begin{table}[htbp]
\centering
\small
\caption{The bootstrap ratchet. Checkpointing begins only when the
qualified pool exists; every parameter only relaxes as the wall of time
fills.}
\label{tab:ratchet}
\begin{tabularx}{\textwidth}{@{}L{1.9cm}L{4.3cm}L{3.6cm}Y@{}}
\toprule
Eras & Finality & Candidacy tenure & Committee, quorum \\
\midrule
0 to 2 & proof-of-work depth: 10 blocks, 20 seconds & no candidacy exists & no Ring, no checkpoints \\
3 to 9 & Ring checkpoints, full force & at least $\min(15, e - 2)$ eras & $n$ qualified, quorum 45 of 55: tolerates 10 faulters \\
10 to 17 & Ring checkpoints, full force & at least $\min(15, e - 2)$ eras & quorum 41 of 55: tolerates 14 faulters \\
18 on & steady state (R1 complete) & 15 eras, incident-free & quorum 37 of 55: tolerates 18 faulters \\
\bottomrule
\end{tabularx}
\end{table}

The mechanics, stated so an implementer needs no interpretation. During
eras 0 to 2 the network runs on classical proof-of-work depth exactly as
it would during a total outage of the committee: payments settle in
twenty seconds, an honest and temporary inferiority, published here
rather than hidden. At era 3 the first candidacy window opens, under the
tenure rule $\min(15, e - 2)$: a genesis-era registrant with six
irreproachable months may run, because requiring seven and a half years
from a one-era-old network is absurd, while requiring none of anyone
would betray the wall. The quorum is 45 of 55, stricter than the steady
37: with a thin reputation base the committee tolerates ten faulters, not
eighteen, and relaxes only as the population deepens, reaching the full
fault tolerance at era 18 when the first candidates can hold fifteen
eras. If fewer than 55 identities qualify, the Ring is partial: $n$
seats, quorum $\lfloor 2n/3 \rfloor + 1$, activating at all only once at
least 21 qualify; if none qualify, checkpoints simply do not start and
depth finality continues. No seat is ever appointed, held by right, or
inherited; the draw happens every era among the qualified; and no
identity may sit more than three consecutive eras, the fourth draw
excluding it for one era, so the bootstrap class cannot entrench itself
into the steady state. The ratchet has one direction: as time passes,
the wall rises for everyone, including the founders, and the committee
weakens its own supermajority only as the honest population grows. The
bootstrap is the one place this paper asks for patience, and it says so
in numbers: finality degrades to twenty seconds for at most eighteen
months, in exchange for never once trusting an unearned reputation.

\section{Veil and Lumen: Privacy That Can Be Read}\label{sec:veil}

The privacy of ANTUMBRA is private by default and provable on demand:
two faces of one mechanism, not two contradictory modes. The Veil, tails
side, makes balances, amounts and payer-paid links invisible to the
network observer; Section~\ref{sec:tx} described its constructions. The
Lumen, heads side, lets the owner open exactly the chosen share, for the
chosen recipient, for the chosen duration. That is the most profitable
lesson of Zcash, read at first degree: that chain survived regulators who
fled the opaque lineage because its mechanism is understood by looking at
it, shielded pools, viewing keys, explicit disclosure. Regulatory
acceptability is not a question of transparency: it is a question of
readability.

The Lumen operates on three levels. The per-transaction viewing key: the
payer proves one precise payment to its single recipient, revealing
nothing else. The bounded auditor key: an accountant receives a reading
right capped in time and bounded in scope; they see the flows of one
activity, not the identity behind it. The compliance proof: a
non-interactive cryptographic proof establishes a fact, the transferred
amount stays under a ceiling, the committed funds are older than $n$
blocks, the balance covers an engagement, without ever revealing amounts,
addresses or history. Porting those proofs to the Travel Rule and the
source-of-funds declarations is a phase four research work item, and the
paper is deliberately conservative about it: the proof family, succinct
arguments with or without a trusted setup, will be chosen at that date
under the then-current state of the art and an external audit, because a
trusted-setup ceremony is a liability to be justified, not a heritage to
be claimed. The pre-genesis rule stands unchanged: no exotic primitive
before the network exists, and nothing phase four activates without the
same audit discipline as the launch itself.

\begin{table}[htbp]
\centering
\small
\caption{Three schools of privacy and their readability.}
\label{tab:privacy}
\begin{tabularx}{\textwidth}{@{}L{3.1cm}YYY@{}}
\toprule
Mechanism & Monero / NERVA & Zcash & ANTUMBRA \\
\midrule
Default view & private (rings) & transparent; shielding optional & private (rings), no transparent pool \\
Disclosure & binary viewing key & viewing key, selective reveal & three levels: transaction, bounded auditor, proof of fact \\
Proof of source of funds & impossible & possible via zero-knowledge proofs & research item, phase 4, family chosen at that date under audit \\
Data protection posture & neutral & transparent pool engraving relationships & native minimization \\
Regulatory readability & low, assumed opacity & strong & strong: disclosure is a protocol primitive \\
\bottomrule
\end{tabularx}
\end{table}

\subsection{The firewall between identity and money}

One tension of this design must be named and closed, because a reader of
the first version named it correctly: Kleos is a public consensus score,
sponsorship and attestation are on-chain social facts, while the Veil
promises unlinkable money. A reputation money cannot buy, paid for by a
graph anyone can read, sitting next to private payments, is a correlation
surface if the two spheres touch. The firewall is a set of protocol
rules, stated once and enforced everywhere.

First, separation of keys by construction: identity keys and wallet keys
grow from different trees, an identity transaction carries no Veil input
or output and a payment carries no identity field, and the encodings are
disjoint, so a mixed object is rejected as malformed, not merely
discouraged. Second, minimum revelation: an Echo attestation publishes
the witness identity, the target identity, a sign and an era, nothing
else, no memo, no amount, no reference to any output; a sponsorship
publishes two identity keys. Third, and it must be said honestly, the
identity layer is pseudonymous, not anonymous: the graph of who sponsors
whom, who testifies for whom, who sits in the Ring is public, because
auditable reputation is the product; what the firewall protects is the
junction with money, and the wallet side is where the practical
discipline lives, published as implementation requirements: identity
actions and payments never leave the same machine through the same Tor
circuit, and the reference wallet enforces it. The residual exposure is
therefore bounded and stated: correlating an Ember with balances requires
evidence from outside the protocol, never the protocol's own records.
Lumen repairs what Kleos cannot: the score is public, the money stays
dark.

The legal argument fits in one line: a chain masked by default is a chain
that collects no personal data by default. That is the very definition of
data minimization, the foundational principle of the European data
protection regulation: compliance by architecture, not by promise. Where
transparent chains engrave graphs of personal relationships into world
history, ANTUMBRA keeps only encrypted commitments and nullifiers; the
data stays in wallets, under the control of their owners. An authority
can obtain a targeted disclosure; no authority can demand a mass
disclosure, because none exists anywhere. The strategic difference with
the transparent pool of Zcash deserves emphasis: ANTUMBRA maintains
nothing transparent, because disclosure there is provable rather than
structural. One does not choose between a glass district and a bunker:
one lives in the bunker and owns a legally admissible flashlight, lit
only on request.
